\section*{Chapter \ref{ch:s2:mortgage-and-structured-credit-strategies} --- Mortgage and Structured-Credit Strategies}

\begin{solution}{exo:s2:mortgage-and-structured-credit-strategies:1}
$26.05 \times 0.358 = 9.32$ per 100 of face.
\end{solution}

\begin{solution}{exo:s2:mortgage-and-structured-credit-strategies:2}
Under the view $148 - 59 = 88$ basis points (from the unrounded values); without burnout $-326 - 42 = -368$ basis points. The same prices give a
large gain or a large loss depending on the model.
\end{solution}

\begin{solution}{exo:s2:mortgage-and-structured-credit-strategies:3}
The tranche is 3\% of the index, 30 million; $30 \times 0.1169 = 3.51$ million.
\end{solution}

\begin{solution}{exo:s2:mortgage-and-structured-credit-strategies:4}
The IO is paid interest only on the principal still outstanding. When rates fall, borrowers refinance, principal disappears and the interest
with it, so the IO loses. Its value therefore rises with rates, like a short bond position, and hedging a short-bond exposure means buying bonds.
\end{solution}

\begin{solution}{exo:s2:mortgage-and-structured-credit-strategies:5}
Gabaix, Krishnamurthy and Vigneron found that it is priced as if the marginal investor were a specialised MBS arbitrageur, for whom prepayment risk is
not diversified away: what is a wash for the economy is concentrated in the books of the few who hold and hedge mortgages.
\end{solution}

\begin{solution}{exo:s2:mortgage-and-structured-credit-strategies:6}
A senior tranche loses only when many names default together, so its expected loss depends on correlation, which is hard to estimate, not on the
pool's expected loss alone. In the chapter's table the super-senior tranche's expected loss is eighteen times larger at a correlation of 0.45 than
at 0.15 for the same pool; Coval, Jurek and Stafford named this fragility of ratings to modest errors as one cause of the collapse of structured
finance.
\end{solution}

\begin{solution}{exo:s2:mortgage-and-structured-credit-strategies:7}
The IO's price rises from 26.05 to 26.85, the view's gain falls from 35.8\% to 33.5\% and the no-burnout loss grows from 25.5\% to 26.5\%. More
volatility sends more paths into deep refinancing, where the wrong burnout assumption costs most, so the view's edge shrinks against its downside;
a prepayment view is also a view on rate volatility.
\end{solution}

\begin{solution}{exo:s2:mortgage-and-structured-credit-strategies:8}
The OAS is an output of the model: at the same prices, a model without burnout gives the 7.5\% coupon an OAS of $-326$ basis points and the 4.0\% one
42, a difference of $-368$. Nothing is locked in; the 88 basis points are earned only if borrowers behave as the view's model says, and the hedge
ratios depend on the model too.
\end{solution}

\begin{solution}{pb:s2:mortgage-and-structured-credit-strategies:1}
\begin{enumerate}
\item A position in mortgage securities whose value depends on prepayment speeds, taken because the trader's prepayment model differs from the
market's, with rate risk hedged.
\item Turnover on the PSA ramp, a refinancing S-curve with a maximum CPR of 45\%, burnout $e^{-0.25c}$; a pool at 6.5\% paying 6.0\% with 29
years left, on Hull-White paths around a flat 4\% curve, at an OAS of 50 basis points, while new mortgages cost about 5.8\%.
\item The decline of refinancing among borrowers who have already passed up incentives; the view assumes a burnout of 1.0 instead of 0.25.
\item Prepayment risk, a wash in the aggregate, is priced, better explained by MBS-market-wide risk, as if the marginal investor were a specialised
arbitrageur.
\item 26.05 per 100 of face; hedged with 0.51 ten-year Treasuries of face 100, bought.
\item Rates unchanged: 0.0\% (market), $35.8\%$ (view), $16.4\%$ (slower refinancing), $-25.5\%$ (no burnout); rates 100 basis points lower:
10.1\%, 63.6\%, 31.1\% and $-36.0\%$.
\item Falling rates raise the incentive, and the models disagree about what borrowers do with it; rising rates remove it, and all models predict the
same turnover.
\item At the market's prices its OAS rises from 59 basis points on the 4.0\% coupon to 148 on the 7.5\%; without burnout it falls from 42 to
$-326$.
\item Trading tranches against each other, the index or the names, because one's model of correlation or pool losses differs from the prices'.
\item Equity loses less and senior tranches lose more as correlation rises; the pool's expected loss, 3.6\%, does not change.
\item Ratings are extremely fragile to modest errors in the underlying risks, and exposed to systematic risk; many structured products paid off like
economic catastrophe bonds for less compensation than comparable alternatives.
\item At a true correlation of 0.15: mean gain 3.1\% of the tranche, a loss in 11.3\% of pools, wiped out in 5.4\%; at 0.30: $-0.4\%$, 14.5\%,
9.4\%; at 0.45: $-1.4\%$, 15.0\%, 10.9\%.
\item Hedged, the IO gains 35.8\% of its price if the burnout view is right, and loses 25.5\% if borrowers never burn out (63.6\% and $-36.0\%$ after
rates fall 100 basis points).
\item That the market prices borrowers' behaviour as a risk, not as a known function of rates: a single model at one OAS would give a flat line.
\item The hedge is sized for small moves: as rates fall far the IO's value flattens, since little is left to prepay, while the Treasury keeps
rising; as rates rise far the IO keeps gaining while the Treasury's losses slow.
\item Estimate it on cohorts and periods not used to fit it, with the data available at the time, and compare predicted and realised speeds by
coupon, seasoning and loan size.
\item For the model being wrong: the no-burnout loss after a rally, not the expected gain.
\item Default correlation shows up only in rare systematic default waves; a few years without one cannot distinguish 0.15 from 0.30.
\item The OAS relative value across coupons: a cross-sectional ranking by a characteristic with duration hedged.
\item The borrower's option to repay, and the trader's edge is knowing how badly it will be exercised.
\end{enumerate}
\end{solution}

\begin{solution}{iq:s2:mortgage-and-structured-credit-strategies:1}
When rates fall, prepayments speed up and the security's price rises less than a bond's; when rates rise, prepayments slow and its duration extends,
so it falls more: its price is a concave function of rates, because the holder is short the borrowers' refinancing option.
\end{solution}

\begin{solution}{iq:s2:mortgage-and-structured-credit-strategies:2}
Split turnover (seasoning, season of the year, home prices) from refinancing (incentive, loan size, credit score, loan-to-value, burnout, media
effects) and curtailment and defaults; fit on loan-level data by cohort, and test out of sample across refinancing waves.
\end{solution}

\begin{solution}{iq:s2:mortgage-and-structured-credit-strategies:3}
Prepayments speed up and the IO loses; the Treasury hedge gains. How much depends on the speeds, so check realised speeds against the model and
rebalance the hedge as the IO's duration changes, shortening it as the IO loses value.
\end{solution}

\begin{solution}{iq:s2:mortgage-and-structured-credit-strategies:4}
Correlation sensitivity, jump to default of single names, spread-level scenarios for the whole index, recovery assumptions, and the tail of pool
losses from a simulation, beyond the first-order deltas.
\end{solution}

\begin{solution}{iq:s2:mortgage-and-structured-credit-strategies:5}
Share the rate paths across pools (common random numbers, generated once), vectorise the cash flows across pools and paths, use fewer paths with
variance reduction, cache the prepayment model's terms, parallelise across cores, and recompute only pools whose inputs changed.
\end{solution}

\begin{solution}{iq:s2:mortgage-and-structured-credit-strategies:6}
The pool loss given the factor $Z$ is $(1-R)\,\Phi\bigl((\Phi^{-1}(p) - \sqrt\rho Z)/\sqrt{1-\rho}\bigr)$; its expectation over $Z$ is
$(1-R)\,P(\sqrt\rho Z + \sqrt{1-\rho}\,\varepsilon < \Phi^{-1}(p)) = (1-R)p$, whatever $\rho$. A tranche's loss is a nonlinear, capped function of the
pool loss, so its expectation depends on the pool loss's whole distribution, which correlation spreads out.
\end{solution}
